\documentclass[11pt,a4paper]{article}
\usepackage{fontspec}
\usepackage{polyglossia}
\setmainlanguage{russian}
\setotherlanguage{english}
\setmainfont{DejaVu Serif}
\setsansfont{DejaVu Sans}
\setmonofont{DejaVu Sans Mono}
\usepackage{amsmath,amssymb}
\usepackage[left=20mm, top=20mm, bottom=20mm, right=20mm, headheight=14pt, headsep=8pt, footskip=12pt]{geometry}
\usepackage{titlesec}
\usepackage{enumitem}
\usepackage{multirow}
\usepackage{xcolor}
\usepackage{tcolorbox}
\tcbuselibrary{skins,breakable}
\usepackage{hyperref}
\hypersetup{colorlinks=true, linkcolor=blue!50!black, citecolor=blue!50!black, urlcolor=blue!50!black}
\usepackage{fancyhdr}
\usepackage{booktabs}
\usepackage{array}
\usepackage{tocloft}

\renewcommand{\cftsecleader}{\cftdotfill{\cftdotsep}}
\renewcommand{\cftsubsecleader}{\cftdotfill{\cftdotsep}}
\renewcommand{\cftsubsubsecleader}{\cftdotfill{\cftdotsep}}

\renewcommand{\cftsecpagefont}{\normalfont}
\renewcommand{\cftsubsecpagefont}{\normalfont}
\renewcommand{\cftsubsubsecpagefont}{\normalfont}

\setlength{\cftsecnumwidth}{2.5em}
\setlength{\cftsubsecnumwidth}{3.2em}
\setlength{\cftsubsubsecnumwidth}{4.2em}

\setlength{\cftsecindent}{0em}
\setlength{\cftsubsecindent}{1.5em}
\setlength{\cftsubsubsecindent}{3em}

\definecolor{mainblue}{RGB}{30,60,110}
\definecolor{boxbg}{RGB}{237,242,250}
\definecolor{boxframe}{RGB}{140,170,210}

\titleformat{\section}
  {\normalfont\Large\bfseries\color{mainblue}}
  {\thesection}{1em}{}
\titlespacing*{\section}{0pt}{3.5ex plus 1ex minus .2ex}{2.3ex plus .2ex}

\titleformat{\subsection}
  {\normalfont\large\bfseries\color{mainblue}}
  {\thesubsection}{1em}{}
\titlespacing*{\subsection}{0pt}{3ex plus 1ex minus .2ex}{1.5ex plus .2ex}

\titleformat{\subsubsection}
  {\normalfont\normalsize\bfseries\color{mainblue}}
  {\thesubsubsection}{1em}{}
\titlespacing*{\subsubsection}{0pt}{2.5ex plus 1ex minus .2ex}{1ex plus .2ex}

\newtcolorbox{formulabox}{colback=boxbg,colframe=boxframe,boxrule=0.6pt,arc=2mm,left=6pt,right=6pt,top=4pt,bottom=4pt,breakable}

\setlength{\parskip}{0.4em}
\emergencystretch=3em
\setlength{\parindent}{0pt}

\pagestyle{fancy}
\fancyhf{}
\fancyhead[L]{\small\textit{Макроэкономика-2. Блок 4. Тема 1}}
\fancyhead[R]{\small\textit{ИП ФЭН, 2026-2027}}
\fancyfoot[C]{\small\thepage}
\renewcommand{\headrulewidth}{0.3pt}

\begin{document}
\begin{titlepage}
    \thispagestyle{empty}
    \centering

    \textcolor{mainblue}{\rule{\textwidth}{1.2pt}}

    \vspace{1.2cm}

    {\color{mainblue}\Huge\bfseries Макроэкономика-2\par}

    \vspace{0.6cm}

    {\color{mainblue}\Large Блок 4. Проблемы фискальной\\ и монетарной политики\par}

    \vspace{0.6cm}

    {\color{mainblue}\large Тема 1. Макроэкономическая политика:\\ цели и ограничения\par}

    \vspace{0.6cm}

    \textcolor{mainblue}{\rule{\textwidth}{0.4pt}}

    \vfill


    \vspace{1cm}

    \textcolor{mainblue}{\rule{\textwidth}{1.2pt}}
\end{titlepage}

\newpage
\tableofcontents
\newpage

\section{Введение: зачем нужна эта лекция}

Предшествующие блоки курса (модели IS--LM, AD--AS и другие) дают фрагментарные и достаточно абстрактные знания о макроэкономической политике. Задача четвёртого блока --- создать целостное представление о роли фискальной (бюджетной) и монетарной (денежно-кредитной) политики и показать, как теория связана с практикой их построения. Лекция выстроена вокруг трёх идей, которые стоит держать в голове с самого начала.

Первая идея: любое решение правительства и центрального банка сопряжено с \textbf{альтернативными издержками}. Нельзя одновременно достичь всех благих целей: стремление к одной цели почти всегда делает другую труднодостижимой, а ресурсы государства ограничены бюджетным ограничением.

Вторая идея: у макроэкономической политики есть как \textbf{позитивные}, так и \textbf{нормативные} аспекты. Позитивный аспект отвечает на вопрос «как политика устроена и почему она именно такая, какая есть» (например, почему правительства систематически допускают дефицит). Нормативный аспект отвечает на вопрос «какой политика должна быть» (например, какая фискальная политика считается устойчивой).

Третья идея: построение фискальной и монетарной политики зависит от \textbf{институционального устройства} экономики, прежде всего от устройства финансовой системы. Одни и те же инструменты работают по-разному в зависимости от того, открыта ли экономика, фиксирован ли курс, как финансируется бюджет и какой у страны доступ к кредитным рынкам.

Весь четвёртый блок включает восемь тем: макроэкономическая политика (цели и ограничения); политика и инфляция; оптимальность и динамическая согласованность политики; режимы монетарной политики; валютные кризисы; политика в развивающихся странах; политика в развитых странах; политика в России. Эта лекция --- первая, вводная, и она разбирает цели политики, ограничения на неё, бюджетное ограничение правительства и центрального банка и одну из двух ключевых проблем политики --- склонность к бюджетному дефициту.

Лекция отвечает на шесть вопросов. Почему макроэкономическая политика не может достичь всех благих целей одновременно? Как интерпретировать данные о состоянии бюджета, если экономика переживает циклические колебания? Почему бюджетное ограничение правительства --- это одновременно ограничение и на фискальную, и на монетарную политику? Зачем накапливать государственный долг или стабилизационный фонд? Что является причиной хронических бюджетных дефицитов? Почему и когда рост государственного долга создаёт проблемы? Дальнейшее изложение построено так, чтобы последовательно ответить на каждый из них.

\section{Цели фискальной и монетарной политики}

Прежде чем говорить об ограничениях, нужно понять, что именно политика пытается сделать. Цели фискальной политики делятся на микроэкономические и макроэкономические.

\textbf{Микроэкономические цели фискальной политики} --- это решение проблем провалов рынка (когда рынок сам по себе не приводит к эффективному распределению ресурсов, например при наличии общественных благ или внешних эффектов) и перераспределение доходов (налоги и трансферты меняют распределение дохода между группами населения).

\textbf{Макроэкономические цели фискальной политики} --- это краткосрочная стабилизация выпуска, занятости и инфляции (сглаживание колебаний деловой активности с помощью государственных расходов и налогов) и стимулирование долгосрочного экономического роста.

\textbf{Цели монетарной политики} иные по составу: стабильность функционирования финансовой системы; стабильный и низкий уровень инфляции; стабильность динамики валютного курса; краткосрочная стабилизация выпуска и занятости.

Уже здесь видно, откуда берётся конфликт целей. Например, краткосрочная стабилизация выпуска и занятости может требовать смягчения денежной политики, а стабильно низкая инфляция и стабильный курс --- её ужесточения. Фискальная политика, нацеленная на краткосрочную стабилизацию, расходует тот же бюджет, который нужен для стимулирования долгосрочного роста и обслуживания долга. Следующий раздел показывает, какие именно ограничения мешают политике достичь всех целей сразу.

\section{Ограничения на макроэкономическую политику}

Ограничения можно разделить на две большие группы. Первая группа связана с тем, как цели и инструменты политики взаимодействуют друг с другом и с поведением экономических агентов: это взаимосвязанность и противоречивость целей, внутренние и внешние лаги, комплементарность инструментов, институциональное устройство экономики (включая стратегическое взаимодействие с правительством и другими центральными банками), а также рациональные ожидания и доверие к политике. Вторая группа --- это бюджетное ограничение: зависимость бюджета от состояния экономической активности, условные обязательства правительства, ограничение устойчивости (sustainability) и ограничение осуществимости (feasibility). Сначала разберём первую группу, затем перейдём к бюджетному ограничению, которому посвящена основная часть лекции.

\subsection{Взаимосвязанность и противоречивость целей}

\subsubsection{Кривая Филлипса: выбор между выпуском и инфляцией}

Первое проявление противоречивости целей --- компромисс между выпуском и инфляцией (output--inflation trade-off), который описывается кривой Филлипса. На слайде лекции он проиллюстрирован данными по безработице и инфляции в США за 1961--1999 гг.: на графике по горизонтали отложен уровень безработицы (примерно от 3 до 11\%), по вертикали --- уровень инфляции (от нуля до примерно 14\%), и облако точек в целом показывает, что более низкой безработице соответствует более высокая инфляция в рамках одного и того же периода.

Из этого следуют два вывода, различающиеся по горизонту. В \textbf{краткосрочном периоде} политика регулирования совокупного спроса не может одновременно добиться увеличения выпуска (то есть снижения безработицы) и сокращения уровня инфляции: стимулируя спрос, правительство и центральный банк получают рост выпуска ценой более высокой инфляции, а подавляя инфляцию --- ценой потерь выпуска и роста безработицы. В \textbf{долгосрочном периоде} выбора между инфляцией и безработицей не существует: систематически «покупать» более низкую безработицу ценой более высокой инфляции нельзя. Таким образом, даже если политике доступны мощные инструменты спроса, они не позволяют одновременно достичь и максимального выпуска, и нулевой инфляции в краткосрочном периоде, а в долгосрочном периоде сама возможность такого выбора исчезает.

\subsubsection{Инструменты монетарной политики в закрытой экономике}

Чтобы понять, чем именно центральный банк управляет и как это связано с процентной ставкой, полезно рассмотреть рынок денег. На слайде показаны две альтернативные схемы, зависящие от того, что именно регулирует центральный банк.

Если центральный банк регулирует \textbf{денежную массу}, то предложение денег фиксировано на уровне $\bar M/P$ (вертикальная линия на графике в координатах «процентная ставка $i$ --- реальные денежные остатки $M/P$»). Спрос на деньги убывает по ставке процента, потому что чем выше ставка, тем дороже держать деньги, не приносящие процента. Сдвиги спроса на деньги (например, с $m_0^d$ на $m_1^d$ и $m_2^d$) при фиксированном предложении меняют равновесную ставку (на графике --- с $i_0$ на $i_1$ и $i_2$): таким образом, \textbf{спрос на деньги определяет равновесную ставку процента}.

Если же центральный банк регулирует \textbf{ставку процента}, то он поддерживает её на заданном уровне $\bar i$ (горизонтальная линия), а предложение денег становится эндогенным: при сдвигах спроса на деньги (с $m_0^d$ на $m_1^d$ и $m_2^d$) центральный банк вынужден менять денежную массу ($M_0/P$, $M_1/P$, $M_2/P$), чтобы удержать ставку. Следовательно, \textbf{спрос на деньги определяет равновесную денежную массу}. Центральный банк не может одновременно зафиксировать и денежную массу, и ставку процента: выбирая один из этих инструментов, он отдаёт другую переменную рынку.

\subsubsection{Инструменты монетарной политики в открытой экономике: трилемма о невозможности}

В открытой экономике добавляется ограничение, известное как \textbf{трилемма о невозможности} (Impossible Trinity). Страна не может одновременно иметь три вещи: совершенную мобильность капитала, режим фиксированного обменного курса и независимую национальную монетарную политику. Можно выбрать любые две, но не все три.

Механизм объясняется на примере страны с фиксированным валютным курсом и совершенной мобильностью капитала, в которой центральный банк хочет провести стимулирующую монетарную политику. На диаграмме IS--LM--BP (по вертикали ставка процента $r$, по горизонтали выпуск $Y$) равновесие вначале находится в точке 0, которая совпадает с точкой 2 (в ней IS, LM$_0$ и BP пересекаются, а платёжный баланс сбалансирован; линия BP горизонтальна на уровне $r_0$, поскольку капитал совершенно мобилен). Рассмотрим последовательность событий:
\begin{enumerate}[leftmargin=1.6em]
\item Центральный банк увеличивает денежную массу ($M\uparrow$), и кривая LM сдвигается вправо (с LM$_0$ на LM$_1$). Равновесие перемещается в точку 1: ставка падает с $r_0$ до $r_1$, а выпуск растёт с $Y_0$ до $Y_1$.
\item Но более низкая ставка внутри страны означает отток капитала (ставка $i\downarrow$ ведёт к снижению сальдо счёта капитала, $CF\downarrow$), и платёжный баланс становится дефицитным ($BP<0$).
\item Это означает, что спрос нашей страны на зарубежные товары и активы превышает спрос иностранцев на наши товары и активы, а значит спрос на отечественную валюту (на рубль) сокращается и её курс начинает падать.
\item Но курс фиксирован, поэтому центральный банк вынужден поддерживать его, \textbf{сокращая} предложение денег и золотовалютные резервы (ЗВР): $M\downarrow$. Кривая LM возвращается обратно, и экономика снова оказывается в точке 2 (то есть там же, где и начинала).
\end{enumerate}
Итог: при фиксированном курсе и совершенной мобильности капитала центральный банк не может осуществлять независимую стимулирующую политику, потому что любая попытка смягчить условия вызывает такой отток капитала, который заставляет его откатить стимул назад.

\subsubsection{Стабильность или эффективность: кредитор последней инстанции}

Ещё одна форма противоречия целей --- конфликт между стабильностью финансовой системы и эффективностью финансовых посредников. Вопрос состоит в том, должен ли центральный банк активно использовать свою функцию «кредитора последней надежды» (the lender of last resort). Аргумент \textbf{за}: ожидаемое вмешательство центрального банка снижает вероятность масштабной банковской паники (если вкладчики уверены, что банку помогут, у них нет причин спешить забирать деньги). Аргумент \textbf{против}: проблемы финансовых посредников, как правило, свидетельствуют об их неэффективности, а значит, помощь неэффективным институтам поощряет неэффективность (и подталкивает к чрезмерному риску). Ключевое методологическое замечание: необходимо различать проблемы \textbf{ликвидности} (liquidity) и проблемы \textbf{платёжеспособности} (solvency). Временная нехватка ликвидных средств у в целом платёжеспособного банка --- повод для помощи; неплатёжеспособность --- нет.

\subsection{Лаги политики}

Любая политика действует не мгновенно. Лекция различает внутренние и внешние лаги.

\textbf{Внутренние лаги} --- задержки на стороне самих органов, проводящих политику. Они складываются из трёх частей: от момента шока до осознания необходимости вмешательства; от осознания необходимости вмешательства до принятия решения; от принятия решения до применения инструмента политики. Например, шок происходит сегодня, статистика покажет его через несколько недель (первая часть), затем правительство должно согласовать пакет мер (вторая часть) и, наконец, реально изменить налоги или расходы (третья часть).

\textbf{Внешние лаги} --- это задержка от применения инструмента до начала его воздействия на экономику. Они определяются поведением частного сектора и скоростью, с которой изменения ставки, налогов или расходов доходят до решений домохозяйств и фирм.

Следствие: к тому моменту, когда политика начинает действовать, экономика может находиться уже в другой фазе цикла, и политика, задуманная как стабилизирующая, способна усилить колебания.

\subsection{Комплементарность инструментов}

Правительство и центральный банк обладают набором \textbf{комплементарных} инструментов, которые с разной степенью эффективности позволяют воздействовать на одни и те же переменные и добиваться одних и тех же целей. Для достижения заданных целей использование инструментов одной из политик выступает как ограничение на использование инструментов другой политики. Например, если правительство увеличивает дефицит для стимулирования спроса, центральному банку, чтобы удержать инфляцию на заданном уровне, придётся быть жёстче, чем в противном случае. Поэтому комплементарность носит \textbf{стратегический} характер, то есть порождает стратегическое взаимодействие фискальной и монетарной политики.

Более того, правительство и центральный банк могут меняться ролями: говорят о \textbf{квази-монетарной политике правительства} и \textbf{квази-фискальной политике центрального банка} (в лекции эти понятия лишь названы, без подробного разбора). Практический вывод тот же: границы между двумя политиками размыты, и бюджетные и денежные решения нельзя рассматривать независимо.

\subsection{Конфликты интересов правительства и центрального банка}

Когда у правительства и центрального банка есть различающиеся интересы, ключевую роль играет форма их стратегического взаимодействия. Лекция выделяет четыре формы: одновременное и независимое взаимодействие (по Курно), когда каждый выбирает свою политику, не зная выбора другого; фискальное лидерство (по Штакельбергу), когда правительство выбирает политику первым, а центральный банк реагирует на неё; монетарное лидерство (по Штакельбергу), когда первым ходит центральный банк; и координация центрального банка и правительства. Существенное замечание: не всегда различие интересов приводит к снижению общественного благосостояния.

\subsection{Ожидания агентов и доверие к политике}

Политика проводится в мире, где частные агенты формируют ожидания и реагируют на саму политику. Лекция упоминает три связанных результата.

\textbf{Критика Лукаса} (Lucas, 1976): попытки использовать при проведении политики устойчивые статистические взаимосвязи между переменными несостоятельны, потому что сама политика, влияя на ожидания агентов, меняет сложившиеся взаимосвязи. Значит, макроэкономическая политика должна учитывать \textbf{рациональную реакцию} экономических агентов.

\textbf{Неэффективность макроэкономической политики} (Sargent \& Wallace, 1975), или Policy Irrelevance Proposition: любые систематические изменения в макроэкономической политике не воздействуют на реальную экономику. В условиях рациональных ожиданий только \textbf{непредвиденные} изменения политики могут повлиять на реальную экономику в краткосрочной перспективе (то, что агенты предвидят, они учитывают заранее).

\textbf{Динамическая несогласованность политики} (Kydland \& Prescott, 1977) --- третий результат. В этой лекции он лишь назван; подробно он разбирается в отдельной теме блока («Оптимальность и динамическая согласованность политики»).

\section{Бюджетное ограничение}

\subsection{Бюджет и экономическая активность: эмпирические факты}

Бюджет не является набором параметров, которые правительство выбирает независимо от экономики: он сам зависит от состояния экономики. Лекция начинает с нескольких эмпирических иллюстраций.

\textbf{Налоги и расходы в зависимости от уровня развития.} На двух диаграммах рассеяния по оси абсцисс отложен реальный ВВП на душу населения (в долларах США 1987 г., от нуля до 20\,000), по оси ординат --- налоговые поступления в процентах ВВП (до 50\%) и государственные расходы в процентах ВВП (до 70\%). Точки образуют облако с положительным наклоном: в более богатых странах и налоги, и государственные расходы в среднем составляют более высокую долю ВВП, хотя разброс велик (точки на диаграммах двух цветов; расшифровка групп на слайде не приводится, но видно, что при высоком доходе на душу доля налогов и расходов у разных стран различается очень сильно). Вывод для интерпретации политики: размер государства сам по себе --- не показатель «хорошей» или «плохой» политики, он связан с уровнем развития.

\textbf{Сальдо бюджета стран еврозоны, 1990--2007.} Лекция приводит таблицу сальдо консолидированного бюджета сектора государственного управления в процентах ВВП. Приведём её целиком (по двум слайдам).

\begin{center}
\footnotesize
\setlength{\tabcolsep}{3.2pt}
\begin{tabular}{lrrrrrrrrrrr}
\toprule
\textbf{Страна} & \textbf{90--98} & \textbf{99--07} & \textbf{1999} & \textbf{2000} & \textbf{2001} & \textbf{2002} & \textbf{2003} & \textbf{2004} & \textbf{2005} & \textbf{2006} & \textbf{2007}\\
\midrule
Бельгия & −5 & −0,2 & −0,5 & 0,1 & 0,4 & 0 & 0 & 0 & −2,3 & 0,3 & −0,2\\
Германия & --- & −2,4 & −1,5 & −1,1 & −2,8 & −3,7 & −4 & −3,8 & −3,4 & −1,6 & 0\\
Ирландия & --- & 1,6 & 2,7 & 4,7 & 0,9 & −0,6 & 0,4 & 1,4 & 1,6 & 3 & 0,3\\
Греция & --- & −4,5 & −3,4 & −3,7 & −4,9 & −4,7 & −5,6 & −7,4 & −5,1 & −2,6 & −2,8\\
Испания & --- & 0,1 & −1,4 & −1,1 & −0,6 & −0,5 & −0,2 & −0,3 & 1 & 1,8 & 2,2\\
Франция & --- & −2,6 & −1,8 & −1,5 & −1,6 & −3,2 & −4,1 & −3,6 & −2,9 & −2,4 & −2,7\\
Италия & --- & −2,9 & −1,7 & −2 & −3,1 & −2,9 & −3,5 & −3,5 & −4,2 & −3,4 & −1,9\\
Кипр & --- & −2,7 & −4,3 & −2,3 & −2,2 & −4,4 & −6,5 & −4,4 & −2,4 & −1,2 & 3,3\\
Люксембург & --- & 2,3 & 3,4 & 6 & 6,1 & 2,1 & 0,5 & −1,2 & −0,1 & 1,3 & 2,9\\
Мальта & --- & −5,3 & −7,7 & −6,2 & −6,4 & −5,5 & −9,8 & −4,6 & −3,2 & −2,5 & −1,8\\
Нидерланды & −3 & −0,5 & 0,4 & 1,3 & −0,2 & −2 & −3,1 & −1,7 & −0,3 & 0,5 & 0,4\\
Австрия & --- & −1,5 & −2,2 & −2,1 & 0 & −0,6 & −1,4 & −3,7 & −1,5 & −1,5 & −0,5\\
Португалия & --- & −3,6 & −2,8 & −3,2 & −4,3 & −2,9 & −2,9 & −3,4 & −6,1 & −3,9 & −2,6\\
Словения & --- & −2,4 & −3,1 & −3,8 & −4,5 & −2,5 & −2,7 & −2,3 & −1,5 & −1,2 & −0,1\\
Финляндия & --- & 3,9 & 1,6 & 6,9 & 5 & 4,1 & 2,6 & 2,4 & 2,9 & 4,1 & 5,3\\
\midrule
Еврозона & --- & −1,9 & −1,4 & −1 & −1,9 & −2,5 & −3,1 & −2,9 & −2,5 & −1,3 & −0,6\\
\bottomrule
\end{tabular}

{\footnotesize Источник: European Commission, NCBs \& ECB calculations. Столбцы «90--98» и «99--07» --- средние за периоды.}
\end{center}

Таблица показывает, что внутри одной валютной зоны фискальные позиции сильно различаются: Люксембург и Финляндия в среднем за 1999--2007 гг. имели профицит (2,3\% и 3,9\% ВВП), а Мальта, Греция, Португалия, Италия, Кипр и Франция --- устойчивые дефициты (от −2,6\% до −5,3\% ВВП). Среднее по еврозоне составило −1,9\% ВВП. Это иллюстрирует, что дефициты --- не случайность отдельного года, а устойчивая черта политики в ряде стран.

\textbf{История США.} Два графика охватывают период с 1792 г. по 2017 г. Федеральный дефицит в процентах ВВП на протяжении большей части XIX в. близок к нулю, а иногда отрицателен (то есть профицит); он резко возрастает в годы войн: заметный всплеск около гражданской войны (1860-е гг., порядка 8\%), в период Первой мировой войны и, особенно, во время Второй мировой войны, когда дефицит превышал 25\% ВВП. После войны дефицит остаётся ниже, но в целом более высок и устойчив, чем в XIX в. Валовой государственный долг в процентах ВВП демонстрирует ту же картину: ступеньки при войнах (пик около 120\% ВВП после Второй мировой войны), последующее снижение до минимума (порядка 30--35\% в 1970-х гг.), затем рост в 1980-х гг. и особенно после 2008 г. --- до уровня порядка 120--130\% к 2017--2020 гг. Таким образом, государственный долг в США исторически накапливался в основном в чрезвычайных обстоятельствах и затем постепенно «выгорал» за счёт роста экономики.

\subsection{Как интерпретировать статистику бюджетного дефицита}

Изменения в государственных расходах и налогах отражают не только решения правительства. Они определяются сразу несколькими группами причин: дискреционными изменениями в фискальной политике (так называемой \textbf{фискальной позицией}, fiscal stance), изменениями в уровне деловой активности (деловые циклы), реальными эффектами инфляции, воздействием внешнеэкономических факторов (например, цен на природные ресурсы) и структурными изменениями в экономике (отраслевые сдвиги, демография). Поэтому, глядя на фактический дефицит, нельзя сказать, стала ли политика мягче или жёстче: дефицит может вырасти просто потому, что экономика вошла в рецессию, налоговые поступления упали, а пособия по безработице выросли --- без каких-либо решений правительства.

Для того чтобы охарактеризовать фискальную политику (то есть оценить fiscal stance), необходимо, как минимум, «очистить» бюджетную статистику от \textbf{циклической составляющей}. Это достигается тремя связанными приёмами, которые на слайдах лекции применены к данным разных стран.

\subsubsection{Циклически скорректированный и первичный баланс}

Сальдо бюджета раскладывается на циклическую и структурную (циклически скорректированную) части. Циклическая часть отражает автоматическую реакцию бюджета на разрыв выпуска (разницу между фактическим и потенциальным ВВП): в рецессию налоги падают, а пособия растут, и наоборот в подъём. Эти автоматические изменения называют действием \textbf{автоматических стабилизаторов}. Циклически скорректированное сальдо показывает, каким был бы баланс бюджета при потенциальном уровне выпуска. Кроме того, из баланса исключают расходы на обслуживание долга, получая \textbf{первичный баланс} (первичный дефицит или профицит): он показывает бюджетную позицию без учёта процентов по накопленному долгу, которые определяются прошлыми решениями. \textbf{Циклически скорректированный первичный баланс} (cyclically adjusted primary balance) --- наилучшая из этих характеристик для оценки дискреционной позиции (fiscal stance).

\subsubsection{Фискальный импульс}

\textbf{Фискальный импульс} (fiscal impulse) измеряет, насколько изменилась позиция политики за год: это изменение сальдо бюджета (в процентных пунктах ВВП) с обратным знаком. Если баланс ухудшился (дефицит вырос), импульс положителен --- политика стала более стимулирующей. Импульс, в свою очередь, раскладывается на вклад автоматических стабилизаторов (циклической компоненты), вклад изменения циклически скорректированного первичного баланса (собственно дискреционной позиции, fiscal stance) и вклад изменения процентных расходов.

На слайдах это показано на примере еврозоны за 2008--2010 гг. (проценты ВВП):

\begin{center}
\small
\begin{tabular}{>{\raggedright\arraybackslash}p{8.3cm}rrr}
\toprule
\textbf{Показатель} & \textbf{2008} & \textbf{2009} & \textbf{2010}\\
\midrule
\multicolumn{4}{l}{\textit{Фискальная позиция (в \% ВВП)}}\\
Сальдо бюджета сектора госуправления & −2 & −6,4 & −6,9\\
Циклическая компонента сальдо & 0,9 & −1,4 & −1,4\\
Циклически скорректированное сальдо & −2,9 & −5 & −5,4\\
Расходы на проценты & 3 & 3 & 3,2\\
Циклически скорректированный первичный баланс & 0,1 & −2 & −2,2\\
Антикризисные фискальные пакеты & --- & 1,1 & 0,8\\
\midrule
\multicolumn{4}{l}{\textit{Фискальный импульс (годовые изменения, п.\,п. ВВП)}}\\
Изменение сальдо бюджета & −1,4 & −4,4 & −0,5\\
Фискальный импульс & 1,4 & 4,4 & 0,5\\
\quad циклическая компонента (автостабилизаторы) & 0,3 & 2,4 & 0\\
\quad циклически скорректированный первичный баланс (fiscal stance) & 1 & 2,1 & 0,2\\
\quad расходы на проценты & 0,1 & 0 & 0,2\\
Изменение объёма антикризисных пакетов & --- & 1,1 & −0,3\\
\bottomrule
\end{tabular}

{\footnotesize Источник: European Commission (2009b \& 2009f), ECB calculations.}
\end{center}

Эту таблицу можно проверить арифметикой. Циклически скорректированное сальдо равно фактическому минус циклическая компонента: в 2008 г. $-2-0{,}9=-2{,}9$; в 2009 г. $-6{,}4-(-1{,}4)=-5{,}0$. Циклически скорректированный первичный баланс равен циклически скорректированному сальдо плюс процентные расходы: в 2008 г. $-2{,}9+3=0{,}1$; в 2009 г. $-5+3=-2$. Фискальный импульс равен минус изменению сальдо: в 2009 г. сальдо упало с $-2$ до $-6{,}4$, то есть изменение составило $-4{,}4$, а импульс $+4{,}4$ п.\,п. ВВП. Эти $4{,}4$ п.\,п. складываются почти поровну из автоматической реакции бюджета на рецессию (2,4 п.\,п.) и дискреционного расширения (2,1 п.\,п.); вклад процентных расходов в 2009 г. нулевой. В 2008 г. картина иная: из 1,4 п.\,п. импульса автостабилизаторы дали лишь 0,3 п.\,п., а основная часть (1 п.\,п.) пришлась на дискреционную позицию. (Небольшие расхождения последнего знака в таблице --- результат округления.) Вывод: масштабное увеличение дефицита еврозоны в 2009 г. только наполовину было результатом решений правительств, остальное --- автоматическим следствием рецессии.

\subsubsection{Примеры оценки фискальной позиции: Великобритания, G-20, Россия}

\textbf{G-20.} На слайде показана эволюция фискальных балансов в процентах потенциального ВВП за 2000--2014 гг. отдельно для развитых и развивающихся стран G-20: общий баланс (overall balance), первичный баланс (primary balance) и структурный первичный баланс (structural primary balance). Во всех случаях видно резкое ухудшение балансов в 2008--2009 гг. (в развитых странах общий баланс падает примерно до $-8$\%), за которым следует постепенное восстановление; структурный первичный баланс ухудшается меньше, чем общий, что подтверждает: значительная часть дефицита имела циклический характер.

\textbf{Великобритания.} Министерство финансов Великобритании в своих отчётах раскладывает общий фискальный эффект на несколько частей. График за 1977--2012 гг. показывает три величины: эффект автоматических стабилизаторов, фискальную позицию (fiscal stance) и разрыв выпуска (output gap). Разрыв выпуска резко уходит в минус в начале 1980-х гг. (порядка $-6$\%) и в начале 1990-х гг. (около $-3$\%); за пределами 2007 г. на графике показаны прогнозы. Таблица (Budget 2008) показывает, как общий фискальный эффект получается из составляющих (проценты ВВП, разница между бюджетами 2007 и 2008 гг.):

\begin{center}
\footnotesize
\setlength{\tabcolsep}{3.5pt}
\begin{tabular}{>{\raggedright\arraybackslash}p{4.9cm}rrrrrr}
\toprule
 & \textbf{2006-07} & \textbf{2007-08} & \textbf{2008-09} & \textbf{2009-10} & \textbf{2010-11} & \textbf{2011-12}\\
 & \textit{факт} & \textit{оценка} & \multicolumn{4}{c}{\textit{прогноз}}\\
\midrule
Изменения бюджета (Budget Changes) & 0 & 0 & 0 & −0,1 & −0,1 & −0,2\\
$+$ Недискреционные факторы & −0,3 & 0,3 & 0,7 & 0,4 & 0,4 & 0,3\\
$=$ Изменение фискальной позиции & −0,3 & 0,4 & 0,7 & 0,4 & 0,2 & 0,1\\
$+$ Автоматические стабилизаторы & −0,1 & −0,2 & 0,2 & 0,3 & 0,2 & 0,1\\
$=$ Общий фискальный эффект & −0,4 & 0,2 & 0,9 & 0,7 & 0,4 & 0,2\\
\bottomrule
\end{tabular}

{\footnotesize Источник: Economic \& Fiscal Strategy Report, March 2008, HM Treasury. Расхождения в последнем знаке вызваны округлением.}
\end{center}

Логика таблицы: сначала складываются дискреционные изменения бюджета и недискреционные факторы, получается изменение фискальной позиции; затем к нему добавляется эффект автоматических стабилизаторов, и получается полный фискальный эффект (например, в 2008--09 гг.: $0{,}7+0{,}2=0{,}9$ п.\,п. ВВП). Второй график, охватывающий 2008--2015 гг., показывает те же три компоненты (автостабилизаторы, фискальная позиция, разрыв выпуска) на более позднем отрезке.

\textbf{Сравнение стран по данным МВФ.} Для 2008--2010 гг. на графике показано среднегодовое воздействие дискреционных мер и автоматических стабилизаторов в процентах ВВП (Италия, Германия, Канада, Франция, Япония, Великобритания, США). Примерные соотношения таковы: в Италии, Франции, Японии и Великобритании основную часть общего эффекта обеспечили автоматические стабилизаторы, а в США, Германии и Канаде заметно больше дискреционных мер; наибольший суммарный эффект наблюдается в США. Общий вывод: страны с развитыми системами социальных гарантий опираются больше на автоматические стабилизаторы, а страны с меньшими автоматическими механизмами используют дискреционные меры.

\textbf{Россия.} По данным Васильевой, Власова, Пономаренко («Анализ стабилизационной функции и устойчивости государственных финансов Российской Федерации», Экономический журнал ВШЭ, 2009, т.\,13, №~3) на графике за 2003--2008 гг. сопоставлены фактический профицит бюджета, структурный профицит, структурный и фактический фискальные импульсы и разрыв выпуска. Фактический профицит вырос с примерно 1,8\% ВВП в 2003 г. до 5\% в 2004 г., около 8\% в 2005 г. и 8,5\% в 2006 г., затем снизился примерно до 6\% в 2007--2008 гг.; структурный профицит (очищенный от циклической составляющей) ведёт себя похоже (около 2,5\% в 2003 г., 5,8\% в 2004 г., 8,5\% в 2005 г., 7,9\% в 2006 г., около 5\% в 2007--2008 гг.). Разрыв выпуска был небольшим отрицательным в 2004--2005 гг. (порядка $-1$\%) и стал положительным и растущим: около 1,7\% в 2006 г. и 4\% в 2007 г. Фискальный импульс (и структурный, и фактический) был положительным в 2004--2005 гг. (около 3 п.\,п. ВВП, то есть политика была стимулирующей), в 2006 г. стал близок к нулю, а в 2007 г. заметно отрицательным (порядка $-3$ п.\,п. для структурного и $-2{,}3$ п.\,п. для фактического импульса). Это означает, что в период положительного разрыва выпуска (перегрева) бюджетная политика ужесточалась, то есть имела контрциклический характер, что и является целью стабилизационной функции бюджета.

\subsection{Условные обязательства правительства}

Фактический бюджетный дефицит и государственный долг не отражают всех обязательств правительства. Помимо них существуют \textbf{условные обязательства} (contingent liabilities), которые превращаются в реальные расходы только при наступлении определённых событий. Они бывают двух видов.

\textbf{Явные условные обязательства} (contingent explicit liabilities): гарантированный государством долг частных или государственных компаний; гарантии конвертации долгов по фиксированному обменному курсу. Если компания-заёмщик не сможет рассчитаться, обязательство по долгу ляжет на государство.

\textbf{Неявные условные обязательства} (contingent implicit liabilities): ожидания государственной поддержки банков и компаний в условиях кризиса. Формального обязательства нет, но участники рынка ожидают, что правительство спасёт крупные институты, и если спасения не произойдёт, оно может потерять доверие, а если произойдёт --- получит новые расходы.

Эти обязательства важны, потому что они расширяют бюджетное ограничение: фактический долг может быть невелик, но реальные потенциальные расходы бюджета окажутся значительно больше.

\subsection{Бюджетное ограничение правительства и центрального банка}

\subsubsection{Обозначения и исходное уравнение}

Введём обозначения: $G_t$ --- реальные государственные закупки в периоде $t$; $T_t$ --- реальные чистые налоги; $B_t$ --- номинальный государственный долг; $R$ --- номинальная ставка процента; $M_t$ --- денежная база; $P_t$ --- уровень цен. Бюджетное ограничение в номинальных величинах имеет вид:
\begin{formulabox}
\begin{equation}
P_t(G_t-T_t)=\frac{B_{t+1}}{1+R}-B_t+M_{t+1}-M_t \tag{1}
\end{equation}
\end{formulabox}
Прочитаем его по частям. Слева стоит \textbf{первичный дефицит бюджета} в номинальном выражении: расходы без учёта обслуживания долга минус налоги. Справа --- способы его финансирования. Слагаемое $B_{t+1}/(1+R)-B_t$ --- приращение долга с учётом его обслуживания: правительство выпускает новые облигации с номиналом $B_{t+1}$, которые продаются со скидкой по цене $B_{t+1}/(1+R)$, и одновременно погашает старые облигации на сумму $B_t$. Слагаемое $M_{t+1}-M_t$ --- денежная эмиссия (фактически, покупка центральным банком государственных облигаций: правительство получает деньги, а центральный банк --- долговые обязательства). Итак, дефицит финансируется либо за счёт новых долгов, либо за счёт печатания денег.

Это же уравнение можно записать иначе, выделив обслуживание долга:
\begin{formulabox}
\begin{equation}
P_t(G_t-T_t)+\frac{RB_t}{1+R}=\frac{B_{t+1}-B_t}{1+R}+M_{t+1}-M_t \tag{2}
\end{equation}
\end{formulabox}
Эквивалентность (1) и (2) проверяется прибавлением $RB_t/(1+R)$ к обеим частям (1): справа получаем $B_{t+1}/(1+R)-B_t+RB_t/(1+R)=(B_{t+1}-B_t)/(1+R)$. В форме (2) слева стоят первичный дефицит и \textbf{обслуживание долга} $RB_t/(1+R)$; их сумма называется \textbf{операционным дефицитом бюджета}. Справа --- приращение долга и денежная эмиссия.

\subsubsection{Переход к реальным величинам и сеньораж}

Разделим (1) на $P_t$:
$$G_t-T_t=\frac{B_{t+1}}{(1+R)P_t}-\frac{B_t}{P_t}+\frac{M_{t+1}-M_t}{P_t}.$$
Последнее слагаемое $(M_{t+1}-M_t)/P_t$ --- реальный доход правительства от печатания денег, который называется \textbf{сеньоражем}. Обозначим его $S_t=(M_{t+1}-M_t)/P_t$. Первое слагаемое преобразуем, домножив числитель и знаменатель на $P_{t+1}$ и введя реальный долг $b_t=B_t/P_t$:
$$\frac{B_{t+1}}{(1+R)P_t}=\frac{B_{t+1}/P_{t+1}}{(1+R)(P_t/P_{t+1})}=\frac{b_{t+1}}{(1+R)/(1+\pi_t)},$$
где $\pi_t=P_{t+1}/P_t-1$ --- уровень инфляции. Отношение $(1+R)/(1+\pi_t)$ определяет \textbf{реальную} валовую ставку процента $1+r$ (уравнение Фишера: номинальная ставка, очищенная от инфляции). Обозначив реальный первичный дефицит $d_t=G_t-T_t$, получаем динамическое бюджетное ограничение правительства в реальных величинах:
\begin{formulabox}
\begin{equation}
d_t=\frac{b_{t+1}}{1+r}-b_t+S_t \tag{3}
\end{equation}
\end{formulabox}
или, выражая $b_{t+1}$,
\begin{formulabox}
\begin{equation}
b_{t+1}=(b_t+d_t-S_t)(1+r) \tag{4}
\end{equation}
\end{formulabox}
Смысл (4): реальный долг следующего периода --- это долг этого периода плюс первичный дефицит за вычетом сеньоража, причём вся сумма растёт по реальной ставке процента. Сеньораж играет роль «альтернативного налога»: его инфляционный характер состоит в том, что деньги печатаются, а платят за это держатели денег, чья покупательная способность размывается ростом цен.

\subsubsection{Итерирование вперёд и ограничение «без схемы Понци»}

Решим (4) относительно $b_t$: $b_t=b_{t+1}/(1+r)-(d_t-S_t)$. Подставим такое же выражение для $b_{t+1}$:
$$b_t=-(d_t-S_t)-\frac{d_{t+1}-S_{t+1}}{1+r}+\frac{b_{t+2}}{(1+r)^2},$$
и продолжим подстановку до $N$ периодов вперёд:
\begin{formulabox}
\begin{equation}
b_t=-\sum_{\tau=0}^{N-1}\frac{d_{t+\tau}-S_{t+\tau}}{(1+r)^{\tau}}+\frac{b_{t+N}}{(1+r)^{N}} \tag{5}
\end{equation}
\end{formulabox}
Накладывается условие отсутствия «схемы Понци» (NPG, No-Ponzi-Game): приведённая стоимость государственного долга должна стремиться к нулю на бесконечном временном горизонте,
\begin{formulabox}
\begin{equation}
\lim_{N\to\infty}\frac{b_{t+N}}{(1+r)^{N}}=0 \tag{6}
\end{equation}
\end{formulabox}
Смысл условия: правительство не может бесконечно рефинансировать долг новыми долгами так, чтобы долг рос быстрее процентной ставки; кредиторы не будут давать в долг тому, кто заведомо не собирается расплачиваться. Из (5) и (6) при $N\to\infty$ получаем:
\begin{formulabox}
\begin{align}
b_t&=-\sum_{\tau=0}^{\infty}\frac{d_{t+\tau}-S_{t+\tau}}{(1+r)^{\tau}} \tag{7}\\
\sum_{\tau=0}^{\infty}\frac{G_{t+\tau}}{(1+r)^{\tau}}+b_t&=\sum_{\tau=0}^{\infty}\frac{T_{t+\tau}+S_{t+\tau}}{(1+r)^{\tau}} \tag{8}
\end{align}
\end{formulabox}
Уравнение (8) --- \textbf{межвременное бюджетное ограничение правительства}: в каждый момент времени приведённая стоимость будущих государственных расходов вместе с накопленным долгом должна быть равна приведённой стоимости налогов, включая сеньораж. Иными словами, всё, что государство тратит (и должно по долгам), оно рано или поздно должно покрыть налогами или печатанием денег.

\subsubsection{Устойчивая фискальная политика (fiscal solvency)}

Если правительство накопило положительный долг ($b_t>0$), то из (7) следует, что сумма $\sum_{\tau}(d_{t+\tau}-S_{t+\tau})/(1+r)^{\tau}$ отрицательна: будущие первичные профициты (с учётом сеньоража) должны перевесить будущие дефициты. Другими словами, приведённая стоимость будущих доходов, включая сеньораж, должна превышать приведённую стоимость будущих расходов на величину, равную текущему долгу. В будущем правительству потребуется сводить бюджет с профицитом или полагаться на сеньораж. \textbf{Устойчивая фискальная политика} --- это политика, построенная так, чтобы приведённая стоимость будущих излишков бюджета (с учётом сеньоража) служила достаточным обеспечением накопленного государственного долга.

\subsubsection{Ограничение осуществимости (feasibility)}

Условие устойчивости само по себе не гарантирует реалистичности политики. Могут ли будущие излишки бюджета служить обеспечением сколь угодно большого накопленного долга? Нет. В каждый момент времени величина первичного излишка бюджета ограничена сверху: налоговые сборы ограничены (существует пик кривой Лаффера --- ставка налога, выше которой поступления перестают расти), а государственные расходы не могут быть отрицательной величиной (нельзя сократить расходы глубже нуля). Помимо формальных пределов существуют и практические сложности: неспособность правительства собрать требуемый объём налогов и отсутствие политической решимости снижать расходы. Поэтому помимо \textbf{ограничения устойчивости} (приведённая стоимость профицитов достаточна) необходимо проверять \textbf{ограничение осуществимости} (требуемые профициты реально достижимы).

\subsubsection{Консолидированное бюджетное ограничение}

Уравнение (8) включает сеньораж --- переменную, которую контролирует центральный банк, поэтому оно по существу является \textbf{консолидированным бюджетным ограничением правительства и центрального банка}. Именно поэтому бюджетное ограничение правительства --- это ограничение и на фискальную, и на монетарную политику: если правительство накапливает долг, а налоги повышать нельзя, нагрузка может перейти на центральный банк в форме требования напечатать деньги. Распределение бремени финансирования бюджетного дефицита (стабилизации государственного долга) между налогами и сеньоражем --- это элемент стратегического взаимодействия правительства и центрального банка, о котором шла речь выше.

\subsection{Способы финансирования дефицита: сеньораж на практике}

Лекция ставит вопросы: какими преимуществами и недостатками обладает каждый способ финансирования дефицита, какой предпочтительнее, как дефицит финансируется на практике в развитых и развивающихся странах, какие проблемы возникают при накоплении долга и как им управлять, каковы реальные эффекты денежной эмиссии (инфляции). Ответ на первый и четвёртый вопросы лежит в плоскости последующих тем курса; здесь приводятся эмпирические данные о роли сеньоража.

Для одних стран сеньораж --- не важный источник финансирования бюджетных дефицитов, для других он, напротив, имеет ключевое значение. На графике сеньораж за 1997--1998 гг. рассчитан как отношение прироста денежной базы к номинальному ВВП (источник --- IMF, International Financial Statistics): у большинства стран (Австралия, Канада, Франция, Германия, Италия, Япония, Великобритания, США) он составляет от 0,5 до 2\% ВВП, у ряда стран (Бразилия, Греция, Индия, Пакистан, Румыния, Россия) --- около 2--4\%, а у Болгарии достигает примерно 17\%.

Таблица более детального сопоставления сеньоража (в \% ВВП и в \% государственных расходов):

\begin{center}
\footnotesize
\setlength{\tabcolsep}{4pt}
\begin{tabular}{lrr@{\hspace{1.2cm}}lrr}
\toprule
\textbf{Страна} & \textbf{\% ВВП} & \textbf{\% расходов} & \textbf{Страна} & \textbf{\% ВВП} & \textbf{\% расходов}\\
\midrule
Дания & 0,3943 & 1,0512 & Таиланд & 1,0872 & 6,3018\\
США & 0,4295 & 1,9552 & Филиппины & 1,2251 & 8,9611\\
Канада & 0,4371 & 2,0144 & Индонезия & 1,3908 & 6,8590\\
Великобритания & 0,4373 & 1,2800 & Корея & 1,5690 & 9,6979\\
Бельгия & 0,4910 & 1,0187 & Малайзия & 1,5749 & 5,2696\\
Нидерланды & 0,5352 & 1,0184 & Испания & 2,0267 & 7,5584\\
Франция & 0,5520 & 1,3863 & Колумбия & 2,3178 & 17,5651\\
Норвегия & 0,5630 & 1,4598 & Италия & 2,3572 & 7,4229\\
Финляндия & 0,5966 & 2,1217 & Бразилия & 3,0394 & 13,7132\\
Швейцария & 0,6187 & 6,7397 & Греция & 3,1291 & 10,5149\\
Германия & 0,6869 & 2,3539 & Мексика & 3,7207 & 18,9687\\
Австрия & 0,6940 & 1,8944 & Аргентина & 9,7299 & 62,0003\\
Швеция & 0,7234 & 1,9301 & Чили & 10,3001 & 32,5765\\
Япония & 0,9585 & 5,6200 & Израиль & 14,8424 & 22,2781\\
\bottomrule
\end{tabular}
\end{center}

Показательны последние два столбца: если в Дании, Бельгии и Нидерландах сеньораж покрывает около 1\% государственных расходов, то в Колумбии --- около 17,6\%, в Мексике --- около 19\%, в Чили --- около 32,6\%, а в Аргентине --- 62\%, то есть значительную долю бюджета в этих странах фактически оплачивали держатели национальной валюты через инфляцию. Так что в странах с развитыми финансовыми рынками правительства финансируют дефицит почти исключительно долгом и налогами, а в ряде развивающихся стран --- заметной долей эмиссии.

\subsection{Государственный долг в странах еврозоны и потребности в рефинансировании}

\textbf{Долг, первичные расходы и процентные расходы.} Следующая таблица показывает, как различаются долговые позиции стран еврозоны (в \% ВВП).

\begin{center}
\footnotesize
\setlength{\tabcolsep}{4pt}
\begin{tabular}{lrrrrrrr}
\toprule
 & \multicolumn{3}{c}{\textbf{Долг}} & \textbf{Перв.\ расх.} & \multicolumn{3}{c}{\textbf{Расходы на проценты}}\\
\cmidrule(lr){2-4}\cmidrule(lr){6-8}
\textbf{Страна} & 99--07 & 1998 & 2007 & 99--07 & 99--07 & 1998 & 2007\\
\midrule
Бельгия & 98,8 & 117,1 & 84,6 & 44,4 & 5,4 & 7,4 & 3,9\\
Германия & 63,3 & 60,3 & 65 & 44,1 & 2,9 & 3,4 & 2,8\\
Ирландия & 32,4 & 54 & 25,5 & 32,4 & 1,4 & 3,3 & 0,9\\
Греция & 100,5 & 105,8 & 94,5 & 39,6 & 5,5 & 8,5 & 3,9\\
Испания & 49,3 & 64,1 & 36,2 & 36,4 & 2,4 & 4,2 & 1,6\\
Франция & 61,5 & 59,4 & 64,2 & 49,8 & 2,8 & 3,3 & 2,6\\
Италия & 106,9 & 114,9 & 104 & 42,7 & 5,4 & 7,9 & 4,7\\
Кипр & 64 & 58,4 & 59,8 & 37,9 & 3,3 & 3,1 & 3,1\\
Люксембург & 6,4 & 7,4 & 6,8 & 39,6 & 0,2 & 0,4 & 0,2\\
Мальта & 63,8 & 53,4 & 62,6 & 40,4 & 3,5 & 3,2 & 3,3\\
Нидерланды & 51,8 & 65,7 & 45,4 & 43 & 2,8 & 4,7 & 2,2\\
Австрия & 64,1 & 64,3 & 59,1 & 47,7 & 3 & 3,5 & 2,6\\
Португалия & 57,5 & 52,1 & 63,6 & 42,4 & 2,8 & 3,3 & 2,9\\
Словения & 26,8 & 23,1 & 24,1 & 44,6 & 1,9 & 2,3 & 1,4\\
Финляндия & 41,9 & 48,2 & 35,4 & 47,2 & 2 & 3,5 & 1,4\\
\midrule
Еврозона & 69 & 72,8 & 66,4 & 44 & 3,4 & 4,6 & 3\\
\bottomrule
\end{tabular}

{\footnotesize Источник: European Commission, NCBs \& ECB calculations. В слайдах приведены также средние за 1990--1998 гг. и первичные расходы 1998 и 2007 гг.; здесь они опущены.}
\end{center}

Таблица иллюстрирует связь долга и процентных расходов: страны с высоким долгом (Италия, Греция, Бельгия) платят по долгу в среднем более 5\% ВВП в год, тогда как у стран с низким долгом (Люксембург, Ирландия) процентные расходы составляют доли процента. Кроме того, процентные расходы почти везде снижались между 1998 г. и 2007 г. (например, в Италии с 7,9 до 4,7\% ВВП, в Греции с 8,5 до 3,9\%) вследствие снижения ставок при вступлении в валютный союз, даже там, где долг почти не менялся.

\textbf{Потребность в финансировании.} Для государства важен не только размер долга, но и его структура сроков погашения: долг, который подлежит погашению в ближайшие годы, нужно рефинансировать. Лекция приводит таблицы «Maturing Debt \& Budget Deficit», показывающие для каждой страны по годам (2013--2015 гг. для развитых, 2013--2014 гг. для развивающихся) три величины: объём погашаемого долга (maturing debt), бюджетный дефицит (budget deficit) и их сумму --- \textbf{совокупную потребность в финансировании} (total financing need). Величины приведены в процентах ВВП (единицы измерения на слайдах явно не указаны, но по смыслу таблиц и порядку величин это проценты ВВП). Таблица ниже содержит часть стран из презентации.

\begin{center}
\footnotesize
\setlength{\tabcolsep}{3.5pt}
\begin{tabular}{lrrrrrr}
\toprule
 & \multicolumn{3}{c}{\textbf{2013}} & \multicolumn{3}{c}{\textbf{2014}}\\
\cmidrule(lr){2-4}\cmidrule(lr){5-7}
\textbf{Страна} & Погаш. долг & Дефицит & Всего & Погаш. долг & Дефицит & Всего\\
\midrule
\multicolumn{7}{l}{\textit{Развитые страны}}\\
Япония & 49,2 & 9,8 & 59 & 51,6 & 7 & 58,6\\
Италия & 25,3 & 2,6 & 27,8 & 25,8 & 2,3 & 28,2\\
США & 18,6 & 6,5 & 25,2 & 20,3 & 5,4 & 25,8\\
Португалия & 17,5 & 5,5 & 23 & 17,8 & 4 & 21,8\\
Испания & 14,1 & 6,6 & 20,7 & 15,3 & 6,9 & 22,2\\
Греция & 14,9 & 4,6 & 19,5 & 19,2 & 3,4 & 22,5\\
Франция & 13,4 & 3,7 & 17,1 & 14,1 & 3,5 & 17,6\\
Германия & 7,9 & 0,3 & 8,2 & 7,9 & 0,1 & 8,1\\
Норвегия & 4,3 & −12,3 & −8 & 4,3 & −11,1 & −6,8\\
\textit{Среднее по группе} & 17,9 & 4,8 & 22,7 & 19,1 & 3,9 & 23\\
\midrule
\multicolumn{7}{l}{\textit{Развивающиеся страны}}\\
Египет & 26,1 & 11,3 & 37,4 & 25,1 & 8,7 & 33,8\\
Пакистан & 26,2 & 7 & 33,2 & 24,7 & 7,1 & 31,9\\
Бразилия & 15,9 & 1,2 & 17,1 & 15,4 & 1,7 & 17,2\\
Индия & 4,4 & 8,3 & 12,7 & 4,3 & 8,4 & 12,7\\
Китай & 4,1 & 2,1 & 6,2 & 3,3 & 1,8 & 5,1\\
Россия & 1,2 & 0,3 & 1,6 & 1,1 & 1 & 2,1\\
\textit{Среднее по группе} & 6,1 & 2,6 & 8,8 & 5,7 & 2,5 & 8,3\\
\bottomrule
\end{tabular}
\end{center}

Из таблицы видно, что страны с одинаковым дефицитом могут иметь совершенно разную потребность в рефинансировании. У Японии потребность составляет 59\% (из них 49,2\% --- погашаемый долг), а у России --- всего 1,6\%; у Норвегии она отрицательна ($-8$\%), поскольку профицит бюджета (12,3\% ВВП) превышает погашаемый долг. У США дефицит 6,5\% в 2013 г., но потребность в финансировании 25,2\% из-за большого объёма долга, подлежащего погашению (расхождения в последнем знаке --- округление). Практический вывод: размер дефицита --- лишь часть картины, а сроки погашения определяют, насколько страна уязвима к ухудшению условий заимствований. Это же показывает, почему важен выбор сроков долга (см. раздел об управлении долгом).

\subsection{Долгосрочная динамика государственного долга}

Лекция даёт несколько исторических графиков, на которых видно, как долг ведёт себя на длинных интервалах.

\textbf{Великобритания, 1700--2016 гг.} Отношение долга к доходу (ВВП) в Великобритании колебалось в диапазоне от примерно 20\% до более 250\%. Первый длительный подъём (порядка 150--190\%) приходится на наполеоновские войны (рубеж XVIII--XIX вв.), затем наступает продолжительное снижение вплоть до минимума около 25--30\% к началу XX в.; затем два мировых конфликта дают подъёмы, максимальный из которых после Второй мировой войны достигает 250--260\% ВВП; после этого долг снижается примерно до 25--30\% к концу XX в., после кризиса 2008 г. снова растёт до примерно 85\% к 2016 г. Важно, что рост долга в Великобритании каждый раз был связан с войнами или кризисом, а последующие длительные периоды снижения долга как доли дохода показывают: высокий долг сам по себе не необратим.

\textbf{США, 1939--2019 гг.} Отношение валового федерального долга к ВВП достигает около 120\% к 1946 г., затем систематически снижается до минимума около 30\% к середине 1970-х гг., растёт с 1980-х гг. до примерно 60--65\% к середине 1990-х гг., снижается немного, а затем после 2008 г. возрастает до уровня около 100\% и выше, достигая около 125\% к 2019--2020 гг.

\textbf{Процентные платежи в США, 1910--1990 гг.} Отношение номинальных и реальных процентных платежей правительства к ВНП показывает, что номинальные процентные платежи выросли с нуля в начале XX в. до примерно 4\% ВНП к 1990 г., в то время как реальные платежи (за вычетом инфляции) в период 1940-х и 1970-х гг. были отрицательными: в годы высокой инфляции кредиторы правительства получали реально отрицательные проценты, то есть инфляция «съедала» государственный долг. После начала 1980-х гг. реальные процентные платежи резко выросли до положительных величин порядка 2--4\% ВНП.

\textbf{Международное сравнение.} На графике валового долга в процентах номинального ВВП за 1980--2010 гг. для Японии, Франции, Великобритании и США видно, что японский долг растёт практически непрерывно с примерно 50\% в 1980 г. до около 230\% к 2010 г., тогда как у США, Великобритании и Франции он колеблется в диапазоне примерно от 20--40\% до 80--100\% с ускорением роста после 2007 г. Карта мира на следующем слайде классифицирует страны по уровню валового долга в процентах ВВП градациями от 0--10\% до более 200\%: наиболее высокие уровни наблюдаются у Японии и ряда развитых стран, а многие развивающиеся и сырьевые экономики имеют низкие уровни долга.

\subsection{Динамика отношения долга к ВВП}

Для оценки того, создаёт ли долг проблемы, важно не абсолютное значение долга, а его величина относительно размера экономики, потому что именно рост экономики определяет способность общества обслуживать долг. Выведем динамику отношения долга к ВВП. Пусть $y_t$ --- реальный ВВП, $g_y$ --- темп его роста. Разделим (3) на $y_t$:
$$\frac{d_t}{y_t}=\frac{b_{t+1}}{(1+r)y_t}-\frac{b_t}{y_t}+\frac{S_t}{y_t}.$$
Преобразуем первое слагаемое тем же способом, что и ранее (домножив на $y_{t+1}$):
$$\frac{b_{t+1}}{(1+r)y_t}=\frac{b_{t+1}/y_{t+1}}{(1+r)(y_t/y_{t+1})}=\frac{b_{t+1}/y_{t+1}}{(1+r)/(1+g_y)}.$$
Введём отношения к ВВП: $dy_t$ --- дефицит, $by_t$ --- долг, $Sy_t$ --- сеньораж, каждое в доле ВВП периода $t$. Получаем:
\begin{formulabox}
\begin{equation}
dy_t=\frac{by_{t+1}}{(1+r)/(1+g_y)}-by_t+Sy_t \tag{9}
\end{equation}
\end{formulabox}
Это то же самое уравнение (3), в котором реальная ставка $r$ заменена на скорректированную на рост ставку $(1+r)/(1+g_y)$. Решая относительно $by_{t+1}$ (сравните с (4)):
\begin{formulabox}
\begin{equation}
by_{t+1}=(by_t+dy_t-Sy_t)\frac{1+r}{1+g_y}
\end{equation}
\end{formulabox}
Именно соотношение реальной ставки процента и темпа роста экономики определяет динамику долга как доли ВВП. Если ставка по долгу превышает темп роста ($r>g_y$), то анализ динамики отношения долга к ВВП качественно совпадает с анализом динамики самого долга: долг как доля ВВП растёт, если нет достаточных первичных профицитов.

\textbf{Эмпирика: ставка процента и рост.} Лекция приводит таблицу реальной ставки процента $r$ и темпа роста ВВП $g$ (в процентах) для десяти европейских стран за 1961--1995 гг. (источник --- IMF, OECD Economic Outlook).

\begin{center}
\scriptsize
\setlength{\tabcolsep}{2.2pt}
\begin{tabular}{l|rr|rr|rr|rr|rr|rr}
\toprule
 & \multicolumn{2}{c|}{\textbf{1961--70}} & \multicolumn{2}{c|}{\textbf{1971--75}} & \multicolumn{2}{c|}{\textbf{1976--80}} & \multicolumn{2}{c|}{\textbf{1981--85}} & \multicolumn{2}{c|}{\textbf{1986--90}} & \multicolumn{2}{c}{\textbf{1990--95}}\\
\textbf{Страна} & $r$ & $g$ & $r$ & $g$ & $r$ & $g$ & $r$ & $g$ & $r$ & $g$ & $r$ & $g$\\
\midrule
Австрия & 3,7 & 4,5 & 1,2 & 4 & 3,3 & 3,4 & 4 & 1,3 & 5,2 & 3 & 4,1 & 2,1\\
Бельгия & 3,5 & 4,9 & −0,6 & 3,5 & 3,2 & 3,1 & 5,3 & 0,8 & 6,3 & 3 & 5,5 & 1,5\\
Дания & 2,1 & 4,4 & 2,6 & 2 & 4,5 & 2,5 & 7,9 & 2,7 & 6,3 & 1,4 & 6,5 & 2,3\\
Франция & 1,7 & 5,6 & −0,2 & 3,5 & −0,4 & 3,1 & 4,1 & 1,5 & 6,1 & 3,2 & 5,6 & 1,3\\
Германия & 4,3 & 4,4 & 2,7 & 2,1 & 3,1 & 3,4 & 4,5 & 1,2 & 5,4 & 3,4 & 3,7 & 2\\
Ирландия & 2,5 & 4,3 & −1 & 4,5 & −0,1 & 4,6 & 2,7 & 2,8 & 6,9 & 5,7 & 6 & 4,7\\
Италия & 3 & 5,7 & −2,6 & 2,8 & −2,3 & 4,8 & 3,4 & 1,4 & 4,9 & 3 & 7,1 & 1,2\\
Нидерланды & 1,6 & 5,7 & −0,5 & 3,2 & 2,7 & 2,6 & 5 & 1 & 6,3 & 3,3 & 4,7 & 2\\
Швеция & 1,9 & 4,7 & −0,3 & 2,6 & −0,3 & 1,3 & 3,8 & 1,8 & 5,3 & 2,3 & 5 & −0,2\\
Великобритания & 3,1 & 2,8 & −1,6 & 2,1 & −1,1 & 1,8 & 4,7 & 2 & 3,9 & 3,3 & 5,3 & 1,4\\
\bottomrule
\end{tabular}
\end{center}

Данные показывают смену режима. До 1980 г. реальная ставка в большинстве стран была ниже темпа роста (а в 1970-е гг. часто отрицательной), с 1981 г. во всех приведённых странах ставка превышает темп роста, нередко в разы (в Бельгии в 1986--1990 гг. 6,3\% против 3\%, в Дании в 1981--1985 гг. 7,9\% против 2,7\%). Этот сдвиг объясняет, почему именно с 1980-х гг. в развитых странах проблема долга стала острой.

\textbf{Два режима.} Если $r>g_y$, то макроэкономическая политика должна быть устойчивой: в противном случае правительство вовлечено в игру Понци, и отношение государственного долга к ВВП возрастает экспоненциально или даже быстрее. Рост долга не может быть бесконечным: это может привести к кризису доверия или к ситуации, когда экономика не способна абсорбировать новые заимствования правительства. Если же $r<g_y$, то для любой «разумной» макроэкономической политики отношение долга к ВВП стабилизируется с течением времени, и сам объём долга может расти экспоненциально, не создавая проблем для экономики страны: рост ВВП опережает рост долга и «разбавляет» его.

\subsection{Зачем накапливать государственный долг или стабилизационный фонд}

Долг --- это пассив правительства, стабилизационный фонд --- его актив. Динамика фонда (обозначим его размер $f_t$) зеркальна динамике долга: если долг растёт на величину дефицита, то фонд растёт на величину профицита. Формулы для долга (слева) и фонда (справа) отражают одну и ту же логику с противоположным знаком:
\begin{formulabox}
\begin{align*}
&\textbf{Долг:} & &\textbf{Фонд:}\\
b_{t+1}&=(b_t+d_t-S_t)(1+r) & f_{t+1}&=(f_t-d_t+S_t)(1+r)\\
\lim_{N\to\infty}\frac{b_{t+N}}{(1+r)^N}&=0 & \lim_{N\to\infty}\frac{f_{t+N}}{(1+r)^N}&=0\\
b_t&=-\sum_{\tau=0}^{\infty}\frac{d_{t+\tau}-S_{t+\tau}}{(1+r)^\tau} & f_t&=\sum_{\tau=0}^{\infty}\frac{d_{t+\tau}-S_{t+\tau}}{(1+r)^\tau}\\
\sum_{\tau=0}^{\infty}\frac{G_{t+\tau}}{(1+r)^\tau}+b_t&=\sum_{\tau=0}^{\infty}\frac{T_{t+\tau}+S_{t+\tau}}{(1+r)^\tau} & \sum_{\tau=0}^{\infty}\frac{G_{t+\tau}}{(1+r)^\tau}&=f_t+\sum_{\tau=0}^{\infty}\frac{T_{t+\tau}+S_{t+\tau}}{(1+r)^\tau}
\end{align*}
\end{formulabox}
Слева: долг обеспечивается будущими профицитами. Справа: фонд позволяет в будущем тратить больше, чем собирается налогов: приведённая стоимость расходов равна фонду плюс приведённой стоимости налогов и сеньоража.

Лекция называет четыре причины накапливать долг или фонд. Во-первых, \textbf{контрциклическая фискальная политика}: в рецессию дефицит финансируется долгом (или тратится фонд), в подъём --- профицит идёт на погашение долга (или пополнение фонда). Во-вторых, \textbf{межвременное перераспределение ресурсов}: правительство может переносить потребление из одного периода в другой. В-третьих, \textbf{сглаживание налогообложения}: чтобы налоговые ставки не прыгали вслед за колебаниями расходов, разумно распределять временные расходы на много периодов с помощью долга. В-четвёртых, \textbf{стратегическое взаимодействие политических партий} (этот мотив подробно разбирается ниже в разделе о склонности к дефициту).

\subsection{Управление государственным долгом}

У управления долгом две цели: минимизация ожидаемых издержек по обслуживанию долга и минимизация риска его обслуживания. Эти цели нередко конфликтуют: более дешёвый (например, краткосрочный или плавающий) долг оказывается более рискованным. Важных аспектов два: \textbf{структура долга} (composition) и \textbf{сроки погашения} (maturity).

\textbf{Долговые инструменты} классифицируются по нескольким признакам.
\begin{itemize}[leftmargin=1.5em]
\item \textit{Валюта долга}: долг в иностранной валюте; рыночные инструменты в отечественной валюте; нерыночные инструменты в отечественной валюте.
\item \textit{Тип процентной ставки}: инструменты с фиксированной доходностью; индексированные к уровню цен; с плавающей доходностью.
\item \textit{Срок до погашения}: краткосрочные (не более 1 года); среднесрочные (до 5--10 лет); долгосрочные (более 5--10 лет).
\end{itemize}
Кроме того, существуют предлагаемые опции: облигации, которые можно выкупить досрочно; облигации с правом досрочного погашения; конвертируемые и продлеваемые облигации. Нерыночный долг включает займы в иностранной валюте у правительств других стран, займы, полученные от международных финансовых организаций, и сберегательные займы для частных (персональных) инвесторов.

\section{Две проблемы макроэкономической политики: склонность к дефициту и к инфляции}

Лекция выделяет две системные проблемы, которые нужно объяснить: \textbf{склонность к дефициту} (deficit bias) и \textbf{склонность к инфляции} (inflation bias). Необходимо понять, почему фискальная политика систематически приводит к дефициту бюджета и почему монетарная политика систематически приводит к слишком высокой инфляции. Эти вопросы позитивные, а не нормативные: речь не о том, какой должна быть политика, а о том, почему реальная политика оказывается иной. В данной лекции подробно разбирается только склонность к дефициту; склонность к инфляции поставлена как вопрос и разбирается в следующих темах блока.

\section{Склонность к дефициту (Deficit Bias)}

\subsection{Почему стандартная теория не объясняет систематический дефицит}

Идея активной фискальной политики (стабилизация выпуска и сглаживание налогов) не может объяснить систематический характер \emph{неэффективно высокого} дефицита бюджета. Действительно, согласно этой идее, дефицит должен возникать в рецессии, а профицит --- в подъёме, так что в среднем за цикл бюджет был бы сбалансирован. Реальность же показывает хронические дефициты (пример --- сальдо бюджетов еврозоны в таблице выше, где в большинстве стран среднее за 1999--2007 гг. отрицательно). Значит, в основе позитивного анализа должны лежать действующие механизмы формирования государственного бюджета, а их рассмотрение лежит в плоскости \textbf{новой политической экономии}. Лекция выделяет два потенциальных источника неэффективной политики: неполные знания политиков и общества относительно оптимальной макроэкономической политики; стратегическое взаимодействие в ходе политического процесса --- отложенная стабилизация, политические циклы и стратегическое накопление долга.

\subsection{Неполные знания об оптимальной политике}

Выгоды от высоких государственных расходов и низких налогов являются \textbf{прямыми и очевидными}: избиратель сегодня видит новую дорогу или меньший налог. Издержки же --- более низкие будущие расходы и более высокие будущие налоги, которые необходимы для выполнения бюджетного ограничения (8), --- являются \textbf{непрямыми и менее очевидными}. Если общество и политики недооценивают эти будущие издержки (то есть не в полной мере осознают, что за сегодняшний дефицит придётся заплатить позже), то политический процесс будет систематически смещён в сторону дефицита.

\subsection{Стратегическое взаимодействие политических партий}

Здесь лекция называет три механизма: политик, находящийся у власти, может накопить неэффективно высокий долг, чтобы ограничить расходы своего преемника; несогласие по поводу распределения бремени снижения дефицита может привести к замедлению реформы фискальной политики; дефицит может формироваться из-за попыток избранных лидеров сообщить о своих способностях избирателям (политические циклы). Рассмотрим первые два подробнее.

\subsubsection{Стратегическое накопление долга}

Политики понимают, что в будущем политику могут проводить партии, с чьими взглядами они не согласны: последующие политики предпочитают расходовать ресурсы так, как сегодняшние политики считают нежелательным. Это даёт сегодняшним политикам желание \textbf{ограничить расходы будущих политиков}. Поскольку высокий уровень государственного долга заставляет сокращать будущие государственные расходы (долг нужно обслуживать, и это сужает пространство для других трат), у сегодняшних политиков появляется стимул накапливать долг. По сути, долг используется как инструмент «связывания рук» преемнику.

Лекция приводит две классические модели.

\textbf{Persson и Svensson (1989): несогласие по поводу уровня государственных расходов.} Консервативные политики предпочитают низкие расходы, а либеральные --- высокие. Если предпочтение консерваторами низких расходов достаточно сильно, они будут проводить политику, приводящую к дефициту: накапливая долг, консерваторы заранее лишают будущих либералов возможности много тратить. Важная оговорка: модель не даёт объяснения общей склонности к дефициту, потому что взаимодействие между политиками может также приводить и к профициту бюджета (всё зависит от того, кто у власти сейчас и кто придёт после).

\textbf{Alesina и Tabellini (1990): несогласие по поводу структуры государственных расходов.} Здесь спор идёт не о том, много или мало тратить, а на что именно. Если каждый тип политика полагает, что тип расходов, который выберет другой тип политика, нежелателен, то у политиков \emph{обоих} типов есть стимул накапливать долг, чтобы ограничить будущие расходы оппонента.

\subsubsection{Отложенная стабилизация}

Если ни одна из групп интересов не контролирует политику единолично, взаимодействие между политическими силами может привести к неэффективному дефициту. В частности, неэффективный дефицит может сохраняться из-за того, что каждый политик или каждая группа интересов затягивает достижение соглашения о фискальной реформе в надежде, что другие возьмут на себя большую часть бремени (Alesina \& Drazen, 1991). Каждая из сторон может выбрать отсрочку в переговорах, чтобы попытаться получить для себя больший выигрыш. Выбирая статус-кво вместо немедленных реформ, группа сигнализирует, что для неё реформа сопряжена с большими издержками.

В результате выбор отсрочки может улучшить ожидаемый исход для группы за счёт ухудшения общей экономической ситуации. Итогом может быть \textbf{отложенная стабилизация}, даже если существует политика, которая, по общему признанию, увеличит благосостояние всех. Каждая группа откладывает согласие принять большую долю бремени реформ в надежде, что это сделает другая группа. Чем меньше издержки для группы при принятии большей доли бремени, тем быстрее она решит, что выгоды от согласия перевешивают выгоды от отсрочки. Так возникает \textbf{война на истощение} (war of attrition): стороны ждут, кто первым не выдержит, или ведут переговоры (торг). Пока длится ожидание, дефицит и долг продолжают накапливаться.

\section{Итоги лекции: ответы на ключевые вопросы}

\textbf{Почему макроэкономическая политика не может достичь всех благих целей?} Потому что цели противоречивы (кривая Филлипса, трилемма о невозможности, стабильность против эффективности финансовых посредников), потому что действуют лаги, потому что инструменты правительства и центрального банка комплементарны и порождают стратегическое взаимодействие, и потому что рациональные ожидания меняют эффект политики (критика Лукаса, неэффективность политики).

\textbf{Как интерпретировать данные о бюджете при циклических колебаниях?} Нужно очистить бюджетную статистику от циклической составляющей, рассчитать циклически скорректированный первичный баланс и фискальный импульс, разложив изменение сальдо на автоматические стабилизаторы, дискреционную позицию и процентные расходы. Пример еврозоны: из 4,4 п.\,п. фискального импульса 2009 г. примерно 2,4 п.\,п. дали автоматические стабилизаторы и 2,1 п.\,п. --- дискреционные меры.

\textbf{Почему бюджетное ограничение правительства --- это ограничение на фискальную и монетарную политику?} Потому что в межвременное бюджетное ограничение (8) входит сеньораж, контролируемый центральным банком, и оно является консолидированным ограничением правительства и центрального банка: приведённая стоимость расходов и долга должна быть равна приведённой стоимости налогов и сеньоража.

\textbf{Зачем накапливать долг или стабилизационный фонд?} Для контрциклической политики, межвременного перераспределения ресурсов, сглаживания налогов и (в политической перспективе) из-за стратегического взаимодействия партий.

\textbf{Что является причиной хронических дефицитов?} Неполные знания (выгоды расходов очевидны, издержки отложены), стратегическое накопление долга (Persson--Svensson, Alesina--Tabellini) и отложенная стабилизация (Alesina--Drazen).

\textbf{Почему и когда рост долга создаёт проблемы?} Критическое значение имеет соотношение реальной ставки $r$ и темпа роста $g_y$: при $r>g_y$ отношение долга к ВВП растёт экспоненциально, если нет достаточных первичных профицитов (игра Понци, кризис доверия), а при $r<g_y$ оно стабилизируется при любой «разумной» политике. Помимо этого, проблемы создают структура и сроки погашения долга, условные обязательства и ограничение осуществимости (пик кривой Лаффера и невозможность отрицательных расходов).

\section*{Литература}
\addcontentsline{toc}{section}{Литература}

\begin{itemize}[leftmargin=1.5em]
\item Бурда М., Виплош Ч. (1998). Макроэкономика. Европейский текст. СПб, Судостроение. Гл.\,3, 9, 15--17.
\item Blanchard O. (2001). Macroeconomics. 3rd ed. Prentice Hall. Ch.\,9, 17, 22--27 (переведены на русский).
\item Romer D. (2001). Advanced Macroeconomics. 2nd ed. McGraw-Hill. Ch.\,10--11 (переведены на русский).
\end{itemize}

\end{document}